% das Papierformat zuerst
%\documentclass[a4paper, 11pt]{article}

% deutsche Silbentrennung
%\usepackage[ngerman]{babel}

% wegen deutschen Umlauten
%\usepackage[ansinew]{inputenc}

% hier beginnt das Dokument
%\begin{document}


\thispagestyle{empty}

\begin{figure}[t]
 \centering
 \includegraphics[width=0.6\textwidth]{abb/fh_whv_big}
\end{figure}


\begin{verbatim}


\end{verbatim}

\begin{center}
\Large{Fachhochschule Wilhelmshaven}\\
%\Large{- Campus <Name> -}\\
\end{center}


\begin{center}
\Large{Fachbereich Management, Information \& Technologie}
\end{center}
\begin{verbatim}




\end{verbatim}
\begin{center}
\doublespacing
\textbf{\LARGE{\titleDocument}}\\
\singlespacing
\begin{verbatim}

\end{verbatim}
\end{center}
\begin{verbatim}

\end{verbatim}
\begin{center}

\end{center}
\begin{verbatim}

\end{verbatim}
\begin{center}
\textbf{zur Erlangung des akademischen Grades \\ Bachelor of Science}
\end{center}
\begin{verbatim}






\end{verbatim}
\begin{flushleft}
\begin{tabular}{llll}
\textbf{Thema:} & & Prototypische Implementierung einer SAP UI5-Applikation & \\
& & im SAP Umfeld und Analyse eines effizienten Einsatz von & \\
& & UI-Objekten & \\
& & & \\
\textbf{Autor:} & & Nils Lutz& \\
& & Mat. Nr. 6002109 & \\
& & & \\
\textbf{Eingereicht im:} & & Wintersemester 2014/15 &\\
& & & \\
\textbf{1. Betreuer:} & & Prof. Dr. Hergen Pargmann &\\
\textbf{2. Betreuer:} & & Prof. Dr. Harald Schallner &\\
\end{tabular}
\end{flushleft}